\section{CONCLUSIONS}

Finishing this study the following conclusions were drawn-up:

\paragraph{Superior Integrator Performance.} Throughout two conducted experiments, a higher-order symplectic method
of Leapfrog fourth-order was concluded to be the most accurate and effective one out of all tested.
Although this method is slightly slower than Velocity Verlet or Leapfrog, it produces much more accurate long-term
results for energy and angular momentum conservation.
This preserves phase-space geometry of the problem, while arriving at an accurate analytical solution.

Long-term simulations benefit from the shorter compute time, since a project running for 100 seconds over a
million iterations can in reality be 27 hours of compute-time, which isn't ideal on consumer devices.
Yet, separating this into multi-core allows for nine hours of compute over three cores.
Hence why even a 20-second reduction from the previously most accurate method is a significant achievement,
which also allows for both more accurate and quicker results.

\paragraph{Symplecticity Over Local Order.} The comparative results across both experiments reinforce that
long-term qualitative fidelity in Hamiltonian systems is governed primarily by the structural, symplectic
properties of an integrator rather than by its local order of accuracy. Explicit Euler's monotonic drift,
the bounded oscillatory error of Symplectic Euler, Leapfrog, and Velocity Verlet, and the resolution-dependent
instability of RK4 together confirm that a non-symplectic method, however high its formal order, cannot
guarantee conservation over extended integration windows. Leapfrog-4 combines the best of both regimes:
the higher-order accuracy associated with methods like RK4, and the bounded, non-drifting error behaviour
guaranteed by the symplectic structure of Leapfrog. This is precisely why it consistently outperformed RK4
in the final experiment, despite RK4's fourth-order local accuracy.

\paragraph{Angular Momentum as a Discriminating Diagnostic.} Angular momentum conservation proved to be a
particularly sharp diagnostic of integrator quality, separating the symplectic methods (Symplectic Euler,
Leapfrog, Velocity Verlet, and Leapfrog-4), which conserved angular momentum to floating-point precision
regardless of timestep, from Explicit Euler and RK4, whose errors scaled directly with resolution. This
structural distinction, arising from whether or not the discrete torque $\mathbf{r} \times \mathbf{F}(\mathbf{r})$
vanishes identically at each step, offers a robust and computationally cheap criterion for evaluating candidate
integrators independent of energy-based diagnostics.

\paragraph{Hypothesis Outcome.} The initial hypothesis, that Velocity Verlet would offer the best trade-off
between computational efficiency and energy fidelity among the five baseline integrators, was largely supported
by the first experiment. However, the subsequent development of Leapfrog-4 demonstrated that this trade-off
could be improved upon further: a composed higher-order symplectic scheme retains Velocity Verlet's structural
conservation guarantees while substantially reducing energy error at a given resolution, at a computational
cost far below that of RK4.

\paragraph{Future Work.} The limitations identified in Section 8.4 point directly to the next steps for this
research. Extending the comparison to adaptive or variable timestep schemes would clarify how much of RK4's
coarse-resolution instability can be mitigated without integrating unnecessarily finely, while addressing the
non-trivial challenge of preserving symplecticity under step-size control. Generalising the present two-body
analysis to N-body systems, particularly those exhibiting mean-motion resonances or chaotic dynamics, would
test whether the ranking established here persists when no analytical ground-truth trajectory is available.
Finally, incorporating GPU or distributed-memory parallelisation, along with post-Newtonian corrections for
strong-field or close-orbit regimes, would extend the applicability of this work beyond the idealised,
single-pair, Newtonian setting studied here.

\paragraph{Closing Remark.} Taken together, these results support the broader conclusion that the pursuit of
higher-order accuracy in Hamiltonian simulation should not come at the expense of symplectic structure.
Leapfrog-4 stands out as the strongest candidate among the integrators examined in this study, offering
long-term energy and angular momentum conservation competitive with, and in several respects exceeding, that
of RK4, while remaining computationally tractable for extended orbital simulations.